\documentclass[11pt,a4paper]{article}

\usepackage[utf8]{inputenc}
\usepackage[T1]{fontenc}
\usepackage[brazilian,provide=*]{babel}
\usepackage[margin=2.5cm]{geometry}
\usepackage{amsmath,amssymb,mathtools}
\usepackage{booktabs}
\usepackage{listings}
\usepackage{xcolor}
\usepackage{hyperref}

\hypersetup{
  colorlinks=true,
  linkcolor=blue!50!black,
  urlcolor=blue!50!black,
  citecolor=blue!50!black,
}

\lstset{
  basicstyle=\ttfamily\small,
  breaklines=true,
  columns=fullflexible,
  keywordstyle=\color{blue!60!black}\bfseries,
  commentstyle=\color{gray},
  showstringspaces=false,
  frame=single,
  framesep=3pt,
  xleftmargin=1em,
  xrightmargin=1em,
}

\newcommand{\vx}{v_x}
\newcommand{\vy}{v_y}

\title{Estratégia para um Passo de Advecção\\
Semi-Lagrangeana no \texttt{femns}}
\author{Projeto \texttt{femns} -- documento de discussão (não implementado)}
\date{Julho de 2026}

\begin{document}

\maketitle

\begin{abstract}
Este documento é um resumo de estratégia, não uma especificação
final: descreve como um passo de advecção semi-Lagrangeana poderia
substituir o tratamento explícito atual do termo convectivo em
\texttt{femns}, mapeando cada etapa do método para módulos concretos
do projeto. O objetivo é servir de base de discussão antes de
qualquer mudança na numérica do solver.
\end{abstract}

\tableofcontents

% =====================================================================
\section{Contexto: como a advecção é tratada hoje}
% =====================================================================

Em \texttt{solver.time\_step} (\texttt{src/femns/solver.py}), o termo
convectivo linearizado é montado explicitamente a cada passo e jogado
para o lado direito do sistema:

\begin{lstlisting}[language=Python]
vg = torch.mm(vx_diag, Gvx) + torch.mm(vy_diag, Gvy)
b_sup = (1/dt) * torch.mm(M, vx.unsqueeze(1)) - torch.mm(vg, vx.unsqueeze(1))
b_mid = (1/dt) * torch.mm(M, vy.unsqueeze(1)) - torch.mm(vg, vy.unsqueeze(1))
\end{lstlisting}

Ou seja: a matriz do sistema $A$ (\texttt{build\_system\_matrix}) só
contém os blocos de massa/rigidez ($M/dt + K/Re$) e o acoplamento de
pressão -- ela \emph{não} contém advecção. O termo advectivo é
avaliado com o campo de velocidade do passo anterior e movido para
$b$. Isso equivale a tratar a advecção de forma \textbf{explícita}
(tipo Euler explícito) enquanto difusão e pressão continuam
implícitas.

O efeito prático: o passo de tempo \texttt{dt} fica limitado por uma
condição do tipo CFL ($\text{dt} \lesssim h / |v|$) para a parte
advectiva, mesmo a parte difusiva/pressão sendo incondicionalmente
estável. Isso é a motivação natural para semi-Lagrangeano: o método
trata a advecção de forma que não impõe essa restrição de CFL --
permitindo \texttt{dt} maior -- ao custo de introduzir difusão
numérica pela interpolação.

% =====================================================================
\section{Ideia central do método}
% =====================================================================

A equação de quantidade de movimento tem a forma
\[
\frac{D\mathbf{v}}{Dt} = \frac{\partial \mathbf{v}}{\partial t} +
(\mathbf{v}\cdot\nabla)\mathbf{v} = -\nabla p + \frac{1}{Re}\nabla^2
\mathbf{v}.
\]
O método semi-Lagrangeano discretiza a \textbf{derivada material}
diretamente, em vez de discretizar separadamente o termo temporal e o
convectivo:
\[
\frac{D\mathbf{v}}{Dt}\bigg|_{\mathbf{x}_n}^{t^{n+1}} \approx
\frac{\mathbf{v}^{n+1}(\mathbf{x}_n) -
\mathbf{v}^{n}(\mathbf{x}_d)}{dt},
\]
onde $\mathbf{x}_d$ é o \textbf{ponto de partida} (\emph{departure
point}): a posição de onde a partícula de fluido que chega em
$\mathbf{x}_n$ no tempo $t^{n+1}$ partiu no tempo $t^n$, seguindo a
característica
\[
\frac{d\mathbf{x}}{dt} = \mathbf{v}(\mathbf{x}, t), \qquad
\mathbf{x}(t^{n+1}) = \mathbf{x}_n.
\]
Como $\mathbf{v}^n(\mathbf{x}_d)$ não cai sobre um nó da malha, ele
precisa ser \textbf{interpolado} a partir dos valores nodais nos
vértices do elemento que contém $\mathbf{x}_d$. Substituindo isso na
equação, o termo convectivo desaparece explicitamente do sistema: ele
já está embutido na escolha de $\mathbf{x}_d$. O sistema linear passa
a ter a mesma estrutura ($M/dt + K/Re$ no bloco de velocidade,
acoplamento de pressão inalterado), só troca o RHS de $\frac{1}{dt}
M\mathbf{v}^n - vg\,\mathbf{v}^n$ (advecção explícita) para
$\frac{1}{dt} M\mathbf{v}^n_\star$ (velocidade interpolada no pé da
característica).

\subsection{O sistema em forma de blocos: antes e depois}

Vale deixar isso explícito, porque é a mudança inteira em uma única
equação. Hoje (\texttt{build\_system\_matrix} + \texttt{time\_step}
em \texttt{solver.py}), o sistema resolvido a cada passo é

\[
\underbrace{\begin{bmatrix}
M/dt + K/Re & 0 & -\beta G_x \\
0 & M/dt + K/Re & -\beta G_y \\
\beta D_x & \beta D_y & 0
\end{bmatrix}}_{A}
\begin{bmatrix} \vx^{n+1} \\ \vy^{n+1} \\ \tilde p^{n+1} \end{bmatrix}
=
\begin{bmatrix}
\dfrac{1}{dt} M \vx^{n} - vg\,\vx^{n} \\[4pt]
\dfrac{1}{dt} M \vy^{n} - vg\,\vy^{n} \\[4pt]
0
\end{bmatrix},
\]
\[
vg = \operatorname{diag}(\vx^{n})\,G_{vx} + \operatorname{diag}(\vy^{n})\,G_{vy}
\]
(a advecção linearizada, montada a partir de $\vx^n, \vy^n$ e das
matrizes convectivas $G_{vx}, G_{vy}$ de \texttt{assembly.py}, entra
só no RHS -- $A$ nunca teve termo advectivo).

Com semi-Lagrangeano, $A$ \textbf{não muda nada}: os mesmos blocos
$M/dt + K/Re$ e o mesmo acoplamento de pressão. Muda só o RHS, que
troca $\vx^n, \vy^n$ (valores no próprio nó) pelos valores
interpolados no pé da característica, e o termo $vg$ some por
completo -- não só de $A$, mas do RHS também:

\[
\begin{bmatrix}
M/dt + K/Re & 0 & -\beta G_x \\
0 & M/dt + K/Re & -\beta G_y \\
\beta D_x & \beta D_y & 0
\end{bmatrix}
\begin{bmatrix} \vx^{n+1} \\ \vy^{n+1} \\ \tilde p^{n+1} \end{bmatrix}
=
\begin{bmatrix}
\dfrac{1}{dt} M \vx^{n}_\star \\[4pt]
\dfrac{1}{dt} M \vy^{n}_\star \\[4pt]
0
\end{bmatrix},
\]
onde $\vx^n_\star, \vy^n_\star$ são os campos interpolados nos pés das
características (Seção~\ref{sec:estrategia}), calculados fora do
sistema linear, num pré-processamento por nó.

Duas consequências diretas dessa comparação:
\begin{itemize}
  \item \texttt{build\_system\_matrix} não precisaria de nenhuma
    mudança -- é puramente a montagem de $A$, que já é a mesma nos
    dois casos;
  \item as matrizes $G_{vx}, G_{vy}$ (produzidas por
    \texttt{assembly.assemble\_mini}) deixariam de ser usadas no
    fluxo de \texttt{time\_step} assim que a advecção explícita for
    removida -- hoje elas só servem para montar $vg$. \textbf{Decisão}
    (\S\ref{sec:transicao}): as duas abordagens ficam lado a lado
    durante a validação, então $G_{vx}, G_{vy}$ continuam montadas e
    em uso normalmente até a migração ser confirmada -- só remover
    depois.
\end{itemize}

% =====================================================================
\section{Estratégia de implementação proposta}
\label{sec:estrategia}
% =====================================================================

Seguindo a separação de responsabilidades do projeto (ver README --
\texttt{mesh.py}/\texttt{assembly.py}/\texttt{boundary.py} não sabem
de config nem de I/O, cada um faz uma coisa), a proposta é um módulo
novo, \texttt{src/femns/semi\_lagrangian.py}, com lógica pura de
malha + interpolação -- sem saber de \texttt{torch} esparso, sistema
global ou BiCGSTAB. \texttt{solver.py} consome sua saída.

\subsection{Passo 1 -- Backtrace: encontrar o ponto de partida}

Para cada nó $\mathbf{x}_n$ (incluindo os nós de centroide/bolha do
elemento MINI, já que $vx, vy$ têm um valor em cada um), calcular
\[
\mathbf{x}_d = \mathbf{x}_n - dt\,\mathbf{v}(\mathbf{x}_n, t^n).
\]
Essa é a aproximação de Euler (1\textsuperscript{a} ordem).
\textbf{Decisão}: usar Euler explícito nesta primeira implementação
-- mais simples de codificar e testar, com uma única avaliação de
velocidade por nó. Uma variante mais precisa, RK2/ponto médio
(estimar $\mathbf{x}_d$ com Euler, avaliar a velocidade no meio do
caminho e refazer o passo), fica como evolução futura caso a
validação da \S\ref{sec:transicao} mostre erro de posição relevante
-- não é necessária para a primeira versão funcional.

\subsection{Passo 2 -- Localizar o elemento que contém o ponto de partida}

Esta é a parte algoritmicamente mais delicada, e a única que
introduz um algoritmo de natureza diferente de tudo que já existe no
projeto (que é montagem vetorizada + álgebra linear). Vale detalhar.

\paragraph{Por que não busca global.} Testar os \texttt{ne} elementos
para cada nó é $O(\text{npoints} \times \text{ne})$ -- na malha do
Poiseuille já mencionada no README ($n\approx 19$ mil incógnitas)
isso é inviável por passo de tempo. A alternativa padrão em métodos
de características para FEM é a \textbf{busca por caminhada}
(\emph{walking search}): partir do elemento que já se conhece (o
elemento ao qual o próprio nó $\mathbf{x}_n$ pertence) e "andar" na
malha, elemento vizinho por elemento vizinho, na direção de
$\mathbf{x}_d$, usando só conectividade local -- custo $O(k)$ por nó,
onde $k$ é o número de elementos cruzados (tipicamente pequeno,
porque o deslocamento $dt\,|\mathbf{v}|$ costuma ser da ordem de uma
fração do tamanho de elemento $h$).

\paragraph{Coordenadas baricêntricas reaproveitando \texttt{assembly.py}.}
Para um triângulo de vértices $i,j,k$ (ordenados como em
\texttt{IEN}), as coordenadas baricêntricas de um ponto
$\mathbf{x}=(x,y)$ são razões de área com sinal:
\[
\lambda_i(\mathbf{x}) = \frac{\operatorname{Area}(\mathbf{x}, j, k)}{\operatorname{Area}(i,j,k)}, \quad
\lambda_j(\mathbf{x}) = \frac{\operatorname{Area}(i, \mathbf{x}, k)}{\operatorname{Area}(i,j,k)}, \quad
\lambda_k(\mathbf{x}) = \frac{\operatorname{Area}(i, j, \mathbf{x})}{\operatorname{Area}(i,j,k)},
\]
com $\lambda_i+\lambda_j+\lambda_k=1$ (só precisa calcular duas, a
terceira sai por diferença). Usando a área com sinal
$\operatorname{Area}(a,b,c)=\tfrac12\big[(x_b-x_a)(y_c-y_a) -
(x_c-x_a)(y_b-y_a)\big]$, essas razões coincidem com as funções de
forma P1 do elemento, $\lambda_i \equiv N_i$, que são exatamente
lineares nos mesmos coeficientes $b_i, c_i$ que
\texttt{assembly.assemble\_mini} já calcula para montar $K$ e $M$
(\texttt{bi = Y[vj]-Y[vk]}, \texttt{ci = X[vk]-X[vj]}, antes do
\texttt{torch.abs} aplicado à área -- aqui precisamos do sinal, não
do valor absoluto). Como $\lambda_i(x_i,y_i)=1$ e $\lambda_i$ é afim
com gradiente $(b_i,c_i)/(2A)$, dá para escrever, sem precisar do
termo independente $a_i$ da forma padrão $N_i=(a_i+b_ix+c_iy)/(2A)$:
\[
\lambda_i(\mathbf{x}_d) = 1 + \frac{b_i\,(x_d - x_i) + c_i\,(y_d - y_i)}{2A},
\]
e de forma análoga para $\lambda_j, \lambda_k$ trocando $(i,b_i,c_i)$
por $(j,b_j,c_j)$ e $(k,b_k,c_k)$. O ponto de reaproveitar os
coeficientes de \texttt{assembly.py} é só evitar duas implementações
divergentes do mesmo elemento -- vale bater esse cálculo contra a
fórmula de área com sinal do parágrafo anterior num teste antes de
decidir qual usar.

\paragraph{Algoritmo de caminhada.} Por nó, partindo do elemento ao
qual o próprio nó pertence:
\begin{enumerate}
  \item Calcular $(\lambda_0, \lambda_1, \lambda_2)$ de
    $\mathbf{x}_d$ em relação aos 3 vértices do elemento atual (o nó
    de centroide/bolha não entra aqui, só na interpolação do
    Passo~3).
  \item Se $\lambda_0, \lambda_1, \lambda_2 \ge -\varepsilon$ (uma
    tolerância numérica pequena, não $\ge 0$ estrito -- ver nota de
    robustez abaixo), o ponto está dentro: parar.
  \item Caso contrário, seja $m = \arg\min_i \lambda_i$ o vértice de
    coordenada mais negativa. Isso indica que o ponto está do lado
    de fora da face oposta a $m$ -- na convenção de
    \texttt{mesh.montar\_EToE} (\texttt{faces\_locais = [(0,1),
    (1,2), (2,0)]}, face $f$ liga os vértices locais $f$ e
    $(f+1)\bmod 3$), a face oposta ao vértice local $m$ é a face
    $(m+1)\bmod 3$. Seguir para \texttt{EToE[elem\_atual, (m+1)\%3]}.
  \item Se \texttt{EToE} indicar \texttt{-1} (face de contorno), o
    ponto de partida caiu fora do domínio -- tratar como
    \S\ref{sec:fora}, não continuar a caminhada.
  \item Repetir a partir do passo 1 no novo elemento, até localizar
    ou sair do domínio.
\end{enumerate}

\paragraph{Notas de robustez (relevantes antes de codificar).}
\begin{itemize}
  \item \textbf{Tolerância na comparação com zero}: como as
    coordenadas vêm de aritmética de ponto flutuante, um ponto
    exatamente sobre uma aresta compartilhada (comum quando o
    deslocamento por passo é pequeno e $\mathbf{x}_d$ cai perto do
    próprio nó de origem) pode dar $\lambda_i$ ligeiramente negativa
    por erro de arredondamento mesmo estando "dentro" -- daí a
    tolerância $\varepsilon$ no critério de parada, em vez de
    comparação estrita com zero.
  \item \textbf{Limite de iterações}: em malhas muito distorcidas ou
    perto de cantos não-convexos, a heurística "andar na direção da
    coordenada mais negativa" pode, em casos patológicos, não
    convergir monotonicamente (Devillers, Pion \& Teillaud, 2001,
    documentam isso para busca por caminhada em triangulações e
    propõem variantes mais robustas, ex.\ \emph{orientation walk}).
    Na prática, para uma malha razoável (a exemplo da gerada por
    Gmsh no projeto), poucos saltos bastam -- mas vale por um limite
    de iterações (ex.\ proporcional a $\sqrt{\texttt{ne}}$) com um
    \emph{fallback} explícito (busca local nos vizinhos de
    \texttt{NToN}, ou marcar como não localizado e cair no
    tratamento de fora-do-domínio) em vez de deixar o laço rodar
    indefinidamente.
  \item \textbf{\texttt{EToE} ainda não está conectado ao
    pipeline}: hoje \texttt{montar\_EToE} existe e é testada
    (\texttt{tests/test\_mesh.py}), mas \texttt{scripts/run\_simulation.py}
    não a chama -- só chama \texttt{montar\_NToN}. Parte do trabalho
    de implementação é efetivamente construir \texttt{EToE} no setup
    (sobre \texttt{mesh.IEN}, antes do \texttt{elem\_mini} acrescentar
    o nó de centroide, já que faces são definidas pelos 3 vértices),
    do mesmo jeito que \texttt{NToN} já é construído hoje.
\end{itemize}

\subsection{Passo 3 -- Interpolar a velocidade no pé da característica}

Comparado ao Passo 2, este é simples: uma vez que se sabe o elemento
e as coordenadas baricêntricas $(\lambda_i, \lambda_j, \lambda_k)$ de
$\mathbf{x}_d$, resta avaliar o campo $\vx^n, \vy^n$ nesse ponto.

\paragraph{Fórmula (base nodal, P1 corrigido + bolha).} O ponto
importante é que $\vx^n,\vy^n$ \emph{já são} funções do espaço MINI --
a mesma base usada para montar $K, M, G_x, G_y$ em
\texttt{assembly.py} --, então interpolar não é ajustar uma função
nova aos valores nodais, é simplesmente \textbf{avaliar a função que
já existe} nesse ponto, usando sua própria expansão de base.
\textbf{Correção importante} (a versão original deste documento tinha
isso errado -- ver nota no fim desta seção): as funções de forma não
são a base "crua" $\{\lambda_i,\lambda_j,\lambda_k,
27\lambda_i\lambda_j\lambda_k\}$, e sim a base \emph{nodal}
\[
N_i = \lambda_i - 9\lambda_i\lambda_j\lambda_k, \qquad
N_b = 27\,\lambda_i\lambda_j\lambda_k
\]
(e ciclicamente $N_j, N_k$), de modo que
\[
v_\star(\mathbf{x}_d) = N_i\,v_i + N_j\,v_j + N_k\,v_k + N_b\,v_b
= \lambda_i v_i+\lambda_j v_j+\lambda_k v_k + 9\lambda_i\lambda_j\lambda_k\big(3v_b-(v_i+v_j+v_k)\big),
\]
onde $v_i,v_j,v_k$ são os valores nodais nos 3 vértices do elemento e
$v_b$ é o valor no nó de centroide desse elemento (a 4\textsuperscript{a}
coluna de \texttt{IEN}, criada por \texttt{mesh.elem\_mini}). A
correção $9\lambda_i\lambda_j\lambda_k$ é escolhida exatamente para
que $N_i, N_j, N_k$ zerem no centroide ($N_i(\mathbf{x}_c)=1/3-9/27=0$)
e $N_b(\mathbf{x}_c)=1$ -- ou seja, $\{N_i,N_j,N_k,N_b\}$ tem a
propriedade de Lagrange nos \emph{4} nós (não só nos 3 vértices), e
$\sum N_a \equiv 1$ (partição da unidade genuína, ao contrário da base
crua).

\paragraph{Por que incluir a bolha, e o que $v_b$ significa.} Isso
resolve a dúvida deixada em aberto numa versão anterior deste
documento ("interpolar só a parte P1 ou incluir a bolha?"): a
resposta é incluir, e não é uma escolha de projeto -- é uma questão
de consistência. E, ao contrário do que este documento afirmava antes,
$v_b$ \textbf{é} a velocidade física no centroide do elemento --
exatamente como $v_i,v_j,v_k$ são as velocidades físicas nos vértices
-- porque a base é nodal nos 4 pontos, não só nos 3 vértices. Descartar
o termo da bolha na interpolação seria avaliar uma função
\emph{diferente} da que o resto do solver está de fato resolvendo --
inconsistente com a própria discretização. Como $\lambda_i,\lambda_j,
\lambda_k$ já foram calculados no Passo 2, o termo extra é barato
(mais uma multiplicação por nó).

\paragraph{Nota: correção de um erro deste documento.} A versão
original desta seção afirmava que $v_b$ era "o coeficiente modal da
bolha, não uma velocidade física no centroide", e usava a fórmula
crua $v_\star=\lambda_i v_i+\lambda_j v_j+\lambda_k v_k+27\lambda_i
\lambda_j\lambda_k v_b$ (sem a correção $9\lambda_i\lambda_j\lambda_k$
nos termos de vértice). Isso é \textbf{tecnicamente incorreto} para
este projeto: batendo a matriz de rigidez que a fórmula errada
implica contra \texttt{tests/test\_assembly.py::test\_assemble\_mini\_%
triangulo\_referencia} (verificação simbólica, entrada a entrada), só
a base nodal acima reproduz exatamente a matriz $K$ que
\texttt{assembly.assemble\_mini} de fato monta -- a base crua dá uma
matriz diferente. O erro foi descoberto durante a integração em
\texttt{solver.py} (\S\ref{sec:transicao}): com a fórmula errada, a
pressão da simulação real divergia (mudava de sinal) já na 2ª
iteração, porque o DOF de bolha do campo já resolvido tende a ficar
muito próximo da média dos vértices (esperado para um DOF nodal de
escoamento suave) -- a fórmula crua interpretava isso como bolha
"quase igual à média" \emph{somada} à própria média, dobrando o valor
físico no centroide. Com a base nodal corrigida, a diferença entre
\texttt{explicit} e \texttt{semi\_lagrangian} no Poiseuille de teste
caiu para ${\sim}10^{-4}$--$10^{-3}$ relativo, o comportamento
fisicamente esperado.

\paragraph{Higiene numérica.} A tolerância $\varepsilon$ do critério
de parada do Passo 2 permite $\lambda_i$ ligeiramente negativa
($-\varepsilon \le \lambda_i < 0$) para pontos perto de uma aresta.
Antes de usar os $\lambda$'s como pesos aqui, vale recortá-los para
$[0,1]$ (\texttt{torch.clamp}) e -- se a soma não bater exatamente em
$1$ por causa do recorte -- renormalizar, para garantir que a parte
P1 da interpolação seja sempre uma combinação convexa dos 3 valores
nodais (sem overshoot introduzido pelo próprio critério de tolerância,
distinto do overshoot inerente à interpolação, comentado abaixo).

\paragraph{Vetorização.} Ao contrário do Passo 2 (caminhada,
sequencial em número de saltos), este passo é uma soma ponderada
simples -- assim que se tem, para cada nó, o índice do elemento
encontrado e os $(\lambda_i,\lambda_j,\lambda_k)$, a interpolação de
todos os nós de uma vez é um \texttt{gather} dos 4 valores nodais por
elemento (3 vértices + bolha) seguido de produto e soma, no mesmo
estilo vetorizado que \texttt{assembly.py} já usa para montar as
matrizes de elemento. Não há laço nesta etapa.

\paragraph{Ordem de precisão e overshoot.}
\label{sec:interp-ordem}
Interpolação P1 (+ bolha)
é exata para campos que já são combinação linear dessa base, mas
introduz erro de truncamento de $O(h^2)$ para campos mais suaves --
esse é o mecanismo concreto da "difusão numérica" mencionada na
\S 4 (Impacto esperado). Um efeito relacionado, comum em métodos
semi-Lagrangeanos, é \emph{overshoot}: interpolação linear não
garante que $v_\star$ fique entre o mínimo e o máximo dos valores
nodais do elemento em todas as situações combinadas com o termo de
bolha (que não é monotônico). Staniforth \& Côté (1991), a referência
clássica de revisão de esquemas semi-Lagrangeanos, discutem esse
problema e limitadores para mitigá-lo -- não é necessário para uma
primeira versão funcional, mas vale ter em mente se aparecerem picos
não-físicos de velocidade perto de gradientes acentuados (ex.\ camada
limite).

\paragraph{Teste.} Seguindo o padrão de \texttt{tests/}: escolher
valores nodais arbitrários $v_i,v_j,v_k,v_b$ para um triângulo de
referência (os mesmos usados em \texttt{tests/test\_assembly.py}),
calcular $v_\star$ à mão para um ponto interno conhecido (ex.\ o
próprio centroide, onde a fórmula colapsa para $v_b$ diretamente,
propriedade de Lagrange) e comparar com a saída da função de
interpolação -- valida a fórmula isolada do Passo 2 (localização).
Implementado em \texttt{tests/test\_semi\_lagrangian.py} (funções
\texttt{interpolar\_mini}/\texttt{baricentro}).

\subsection{Passo 4 -- Trocar o termo convectivo em \texttt{solver.py}}

Em \texttt{time\_step}, substituir:
\begin{lstlisting}[language=Python]
vg = torch.mm(vx_diag, Gvx) + torch.mm(vy_diag, Gvy)
b_sup = (1/dt) * torch.mm(M, vx.unsqueeze(1)) - torch.mm(vg, vx.unsqueeze(1))
\end{lstlisting}
por algo como:
\begin{lstlisting}[language=Python]
vx_star, vy_star = semi_lagrangian.calculo_sl(...)
b_sup = (1/dt) * torch.mm(M, vx_star.unsqueeze(1))
\end{lstlisting}
A matriz $A$ (\texttt{build\_system\_matrix}) não muda -- ela já não
tinha advecção. O ganho de \texttt{dt} maior sem instabilidade vem
inteiramente dessa troca no RHS.

\subsection{Pontos de partida fora do domínio}
\label{sec:fora}

Perto de contornos de entrada (\emph{inlet}), o pé da característica
pode cair fora da malha -- fisicamente correto (a partícula "veio de
fora", carregando o valor de contorno). \textbf{Decisão}: se a
caminhada sair do domínio, usar o valor de Dirichlet do contorno por
onde saiu (ex.\ \texttt{inlet: \{vx: 1.0, vy: 0.0\}} no config); se o
contorno não definir esse componente ali (ex.\ \texttt{outlet}, que
só define \texttt{p} no exemplo do README), usar o valor atual do
próprio nó de chegada como aproximação -- equivalente a um passo de
Euler explícito só nesse nó, ou seja, cai de volta no comportamento
de hoje exatamente nos casos em que não há dado de contorno melhor
para usar. \texttt{boundary.py} já resolve qual nome de contorno cada
nó tem (\texttt{assign\_boundary\_names}) -- a mesma lógica de
prioridade dá para reaproveitar aqui.

\subsection{Vetorização em GPU / \texttt{torch}}

A caminhada (Passo 2) é sequencial por natureza (número de passos
varia por nó), o que não vetoriza tão bem quanto o resto do projeto
(que é 100\% montagem vetorizada em \texttt{torch}, ver
\texttt{assembly.py}). Estratégia realista: vetorizar \emph{entre}
nós (todos os nós fazem "um passo de caminhada" em paralelo, um
tensor booleano marca quais já terminaram) e iterar até que todos
convirjam -- número de iterações do laço = máximo de saltos de
elemento entre todos os nós, não \texttt{npoints} $\times$
\texttt{ne}. Vale prototipar em CPU/\texttt{numpy} primeiro (mais
fácil de depurar com \texttt{prints}/asserts) e só depois portar para
\texttt{torch} caso o custo por passo de tempo se mostre relevante.

\subsection{Testes}

Seguindo o padrão de \texttt{tests/} (que testa \texttt{assembly.py}
contra valores analíticos de um triângulo de referência, sem
depender de malha \texttt{.msh}): um teste unitário com 1--2
triângulos e um campo de velocidade constante (translação pura) tem
solução fechada para o ponto de partida e para a interpolação --
bom primeiro caso para validar localização + interpolação
isoladamente, antes de plugar em \texttt{time\_step}.

\subsection{Estratégia de transição e validação}
\label{sec:transicao}

\textbf{Decisão}: as duas abordagens de advecção ficam lado a lado
durante a validação, em vez de substituir a atual de uma vez. Um
parâmetro novo no config (ex.\ \texttt{simulation.advection: explicit
| semi\_lagrangian}, junto de \texttt{dt}/\texttt{iterations}/
\texttt{reynolds}) seleciona qual caminho \texttt{time\_step} usa a
cada passo -- \texttt{explicit} continua montando $vg$ com $G_{vx},
G_{vy}$ como hoje, \texttt{semi\_lagrangian} usa
\texttt{semi\_lagrangian.calculo\_sl}. Isso mantém as
matrizes $G_{vx}, G_{vy}$ em uso normalmente até a migração ser
confirmada.

Validação proposta, no mesmo espírito da tabela comparativa que
valida o BiCGSTAB em \texttt{docs/bicgstab\_solver.pdf}: rodar o
Poiseuille de exemplo (\texttt{configs/poiseuille.yaml}, com
\texttt{iterations} reduzido para testar rápido) com \texttt{dt}
pequeno -- onde a advecção explícita atual é estável e portanto serve
de referência -- nas duas abordagens, e comparar posição de nó
(efeito do \texttt{moving\_point}) e campos de velocidade/pressão ao
longo das iterações. Só depois de bater (dentro de uma tolerância
compatível com a ordem de erro esperada da interpolação P1+bolha, ver
\S\ref{sec:interp-ordem}) faz sentido testar \texttt{dt} maior -- onde
as duas deixam de ser diretamente comparáveis, já que aí é onde o
semi-Lagrangeano deveria continuar estável e o explícito não. Remover
$G_{vx}, G_{vy}$/\texttt{vg} e a opção \texttt{explicit} do config só
depois dessa validação, não antes.

% =====================================================================
\section{Impacto esperado e trade-offs}
% =====================================================================

\begin{itemize}
  \item \textbf{A favor}: remove a restrição de CFL da advecção,
    permitindo \texttt{dt} maior (potencialmente relevante já que o
    README aponta \texttt{dt} pequeno como parte do custo por
    iteração).
  \item \textbf{Contra}: interpolação (mesmo linear) introduz difusão
    numérica -- a solução deixa de ser idêntica à advecção Galerkin
    explícita atual, mesmo em \texttt{dt} pequeno onde ambas seriam
    estáveis. Isso é uma mudança de \emph{numérica}, não só de
    performance.
  \item Complexidade nova: busca por elemento (Passo 2) é um
    algoritmo diferente de tudo que existe hoje no projeto (que é só
    montagem vetorizada + solve linear) -- maior superfície para
    bugs sutis (ex.: caminhada que não converge em malhas muito
    distorcidas).
\end{itemize}

% =====================================================================
\section{Referências}
% =====================================================================

\begin{enumerate}
  \item Pironneau, O. (1982). \emph{On the transport-diffusion
    algorithm and its applications to the Navier--Stokes equations}.
    Numerische Mathematik, 38(3), 309--332. -- artigo original que
    introduz o método de características (semi-Lagrangeano) aplicado
    a elementos finitos para Navier--Stokes; base teórica direta para
    a formulação da \S 2.
  \item Donea, J., \& Huerta, A. (2003). \emph{Finite Element Methods
    for Flow Problems}. Wiley. -- capítulo sobre métodos
    Characteristic-Galerkin/semi-Lagrangeanos, com algoritmos práticos
    de localização de ponto em malhas de elementos finitos (base do
    Passo 2) e discussão de erro de interpolação/difusão numérica
    (\S 4).
  \item Devillers, O., Pion, S., \& Teillaud, M. (2001). \emph{Walking
    in a triangulation}. International Journal of Foundations of
    Computer Science, 13(2), 181--199. -- análise da busca por
    caminhada (\emph{walking search}) usada no Passo 2: condições de
    convergência, casos patológicos e variantes mais robustas que a
    heurística "coordenada mais negativa" usada aqui.
  \item Staniforth, A., \& Côté, J. (1991). \emph{Semi-Lagrangian
    integration schemes for atmospheric models -- a review}. Monthly
    Weather Review, 119(9), 2206--2223. -- revisão clássica de
    esquemas semi-Lagrangeanos; discute ordem de precisão da
    interpolação no pé da característica e o problema de overshoot
    citado no Passo 3.
\end{enumerate}

\end{document}
